\documentclass[12pt,a4paper]{article}

\usepackage[utf8]{inputenc}
\usepackage[T1]{fontenc}
\usepackage[spanish,es-nodecimaldot]{babel}
\usepackage{amsmath}
\usepackage{amssymb}
\usepackage{graphicx}
\usepackage{hyperref}
\usepackage[a4paper,margin=2.5cm]{geometry}
\usepackage{float}
\usepackage{caption}
\usepackage{subcaption}
\usepackage{IEEEtrantools}
\usepackage{cite}
\usepackage{array}
\usepackage{longtable}
\usepackage{xcolor}
\usepackage{enumitem}

\hypersetup{
  colorlinks=true,
  linkcolor=blue!50!black,
  citecolor=blue!50!black,
  urlcolor=blue!50!black,
  pdftitle={Documentación de Credit Simulator},
  pdfauthor={Joselyn Redroban}
}

\setlength{\parindent}{0pt}
\setlength{\parskip}{0.45em}
\renewcommand{\arraystretch}{1.2}

\title{\textbf{Credit Simulator}\\
\large Documentación técnica y académica del proyecto}
\author{Joselyn Redroban}
\date{\today}

\begin{document}

\maketitle

\begin{abstract}
Este documento describe el proyecto Credit Simulator, una aplicación web para
simular créditos y consultar sus tablas de amortización. Se documentan el
propósito, la arquitectura de microservicios, el frontend, los endpoints, la
autenticación con JSON Web Tokens (JWT), la persistencia en PostgreSQL, los
métodos de amortización francés y alemán, el flujo de uso y la infraestructura
Docker Compose. También se registran decisiones de diseño, riesgos, ausencia
de pruebas automatizadas visibles y recomendaciones de mejora. La descripción
se basa en el código y los archivos de configuración del repositorio, por lo
que no se atribuyen al sistema capacidades que no estén implementadas.
\end{abstract}

\textbf{Palabras clave:} simulación de crédito, microservicios, amortización,
JWT, PostgreSQL, ASP.NET Core, React, Docker

\begin{otherlanguage}{english}
\begin{abstract}
This document presents the Credit Simulator project, a web application for
credit simulation and amortization-table consultation. It describes the
microservice architecture, frontend, HTTP endpoints, JWT authentication,
PostgreSQL persistence, French and German amortization methods, usage flow,
Docker Compose infrastructure, and configuration. It also records design
decisions, limitations, security risks, and recommendations. The account is
based on the repository implementation and deliberately avoids claiming
features or measured results that are not present in the project.
\end{abstract}
\end{otherlanguage}

\tableofcontents
\newpage

\section{Introducción}

Credit Simulator permite que una persona registrada explore el costo de un
crédito antes de tomar una decisión. El usuario elige un tipo de crédito,
ingresa el monto y el plazo, selecciona el método de amortización y obtiene
una tabla con cuota, interés, abono a capital y saldo. Además, la aplicación
guarda las simulaciones asociadas al usuario y permite consultar el historial y
descargar el resultado en CSV desde el backend o en PDF desde el frontend.

El proyecto es especialmente útil como ejercicio de integración entre una
interfaz web, APIs independientes, persistencia relacional y cálculos
financieros. Sin embargo, debe entenderse como un simulador educativo y no
como un motor bancario: no se observan comisiones, seguros, impuestos,
calendarios de días reales, tasas variables ni validaciones regulatorias.

\section{Propósito y alcance}

\subsection{Objetivo}

El objetivo es ofrecer una experiencia sencilla para:
\begin{enumerate}[label=\arabic*.]
  \item registrar usuarios e iniciar sesión;
  \item consultar un catálogo de tipos de crédito con tasa anual;
  \item calcular y comparar amortización francesa y alemana;
  \item conservar las simulaciones y sus cuotas;
  \item revisar el historial privado y exportar un resultado.
\end{enumerate}

\subsection{Alcance comprobado}

El repositorio contiene tres servicios ASP.NET Core sobre .NET 8, un frontend
React con Vite, scripts SQL de inicialización, Dockerfiles y una composición
Docker. El catálogo incluye cinco registros semilla: Personal, Hipotecario,
Vehicular, Educativo y PyME. El plazo aceptado por la solicitud de simulación
es de 1 a 600 meses y el monto debe ser positivo.

No se encontró una carpeta de pruebas unitarias o de integración ni proyectos
de pruebas en el árbol revisado. El README registra validaciones manuales
realizadas durante el desarrollo, pero esas comprobaciones no sustituyen una
suite reproducible dentro del repositorio. Esta diferencia es importante para
interpretar cualquier afirmación de calidad.

\section{Arquitectura general}

\subsection{Vista lógica}

La solución aplica separación por responsabilidad y por almacenamiento. Auth
Service administra usuarios y tokens; Credit Catalog Service administra los
tipos de crédito; Simulation Service consulta el catálogo, calcula y persiste
las simulaciones. El navegador consume los servicios mediante Axios.

\begin{figure}[H]
\centering
\fbox{\begin{minipage}{0.92\textwidth}
\centering
\textbf{Frontend React/Nginx}\\[0.3em]
\(\downarrow\) autenticación, catálogo y simulaciones por HTTP\\[0.3em]
\begin{tabular}{c@{\qquad}c@{\qquad}c}
\textbf{Auth Service} & \textbf{Credit Catalog Service} &
\textbf{Simulation Service}\\
\(\downarrow\) & \(\downarrow\) & \(\downarrow\)\\
\texttt{auth\_db} & \texttt{credit\_catalog\_db} &
\texttt{simulation\_db}
\end{tabular}\\[0.4em]
Simulation Service \(\longrightarrow\) Credit Catalog Service por HTTP
\end{minipage}}
\caption{Diagrama lógico simplificado, construido sin archivos externos.}
\label{fig:arquitectura}
\end{figure}

Cada servicio posee su propia base de datos. En consecuencia, no existen
claves foráneas cruzadas entre usuarios, catálogo y simulaciones. Simulation
Service guarda el identificador del usuario y toma una copia del nombre y la
tasa del crédito usados en el momento de la simulación. Esta decisión reduce
la dependencia de lecturas posteriores del catálogo, aunque deja la
consistencia referencial entre servicios bajo responsabilidad de la
aplicación.

\subsection{Comunicación}

El frontend usa por defecto los puertos locales 5001, 5002 y 5003 para las
APIs. Dentro de Docker, Simulation Service usa la URL interna
\texttt{http://credit-catalog-service:8080}. El catálogo se consulta mediante
\texttt{CreditCatalogClient}; el servicio de simulación no accede directamente
a las tablas del catálogo.

\section{Componentes del sistema}

\subsection{Auth Service}

Expone registro e inicio de sesión. El correo se normaliza a minúsculas y se
recorta; la contraseña se almacena como hash BCrypt, no como texto plano.
Cuando las credenciales son válidas se genera un JWT con identificador,
nombre y correo del usuario, emisor, audiencia y una expiración configurada de
60 minutos por defecto. La tabla \texttt{users} incluye una restricción única
para el correo.

\subsection{Credit Catalog Service}

Expone el catálogo de forma de solo lectura. La respuesta contiene identificador,
nombre, descripción y tasa anual. La semilla SQL define tasas positivas para
cinco tipos de crédito. El servicio no implementa en el código revisado
operaciones administrativas para crear o modificar el catálogo.

\subsection{Simulation Service}

Es el componente central del caso de uso. Requiere autenticación JWT,
valida la solicitud mediante anotaciones de datos, consulta el tipo de crédito,
selecciona el calculador según el método, genera las cuotas, calcula totales y
guarda la cabecera junto con sus cuotas. Las consultas de historial y detalle
filtran por el identificador del usuario autenticado; por eso una simulación
de otro usuario se presenta como no encontrada.

\subsection{Frontend}

El frontend está construido con React 18, React Router, Axios, Vite y jsPDF.
Incluye páginas de registro, login, simulador, historial y detalle. El
\texttt{AuthContext} conserva el token y los datos del usuario en
\texttt{localStorage}; un interceptor de Axios añade el encabezado
\texttt{Authorization: Bearer <token>}. La ruta privada impide mostrar las
pantallas principales a una sesión no autenticada.

El PDF se construye en el navegador con jsPDF. El CSV se solicita al endpoint
del backend como un archivo binario. Esta división es práctica, pero produce
dos implementaciones de exportación que podrían divergir en formato o
localización.

\section{Endpoints y flujo de datos}

\small
\begin{longtable}{>{\ttfamily}p{0.22\textwidth}p{0.18\textwidth}p{0.48\textwidth}}
\caption{Endpoints observados en los controladores.}\label{tab:endpoints}\\
\hline
\normalfont Método y ruta & \normalfont Servicio & \normalfont Función\\
\hline
\endfirsthead
\hline
\normalfont Método y ruta & \normalfont Servicio & \normalfont Función\\
\hline
\endhead
\textnormal{POST }/api/auth/register & Auth & Registra un usuario y devuelve JWT.\\
\textnormal{POST }/api/auth/login & Auth & Verifica BCrypt y devuelve JWT.\\
\textnormal{GET }/api/credit-types & Catálogo & Devuelve los tipos de crédito.\\
\textnormal{POST }/api/simulations & Simulación & Crea y persiste una simulación autenticada.\\
\textnormal{GET }/api/simulations & Simulación & Lista el historial del usuario autenticado.\\
\textnormal{GET }/api/simulations/\{id\} & Simulación & Obtiene el detalle propio.\\
\textnormal{GET }/api/simulations/\{id\}/export & Simulación & Descarga el reporte CSV propio.\\
\hline
\end{longtable}
\normalsize

El flujo normal es: (1) registro o login; (2) almacenamiento local del token;
(3) carga del catálogo; (4) envío de monto, plazo, tipo y método; (5)
validación y consulta de la tasa; (6) cálculo y persistencia; (7) presentación
de la tabla; y (8) navegación al detalle, historial o exportación. El
frontend muestra la primera cuota como resumen, mientras que el detalle
presenta todas las cuotas.

\section{Modelo de persistencia}

La base \texttt{auth\_db} contiene \texttt{users(id, nombre, email,
password\_hash, created\_at)}. La base \texttt{credit\_catalog\_db} contiene
\texttt{credit\_types(id, nombre, descripcion, tasa\_anual)}. La base
\texttt{simulation\_db} contiene \texttt{simulations} e
\texttt{installments}; la segunda tiene una clave foránea hacia la primera
con borrado en cascada.

\begin{table}[H]
\centering
\caption{Datos históricos conservados en una simulación.}
\label{tab:persistencia}
\begin{tabular}{>{\raggedright\arraybackslash}p{0.28\textwidth}p{0.58\textwidth}}
\hline
\textbf{Campo} & \textbf{Razón de conservación}\\
\hline
\texttt{usuario\_id} & Filtrar historial y propiedad del recurso.\\
\texttt{credit\_type\_id} & Identificar el registro de catálogo usado.\\
\texttt{tipo\_credito} & Conservar el nombre aunque el catálogo cambie.\\
\texttt{tasa\_anual} & Reproducir la tasa aplicada en el pasado.\\
\texttt{monto}, \texttt{plazo\_meses}, \texttt{metodo} & Reproducir los parámetros de entrada.\\
\texttt{total\_intereses}, \texttt{total\_pagado} & Mostrar agregados sin recalcular la tabla.\\
\hline
\end{tabular}
\end{table}

Los índices declarados para usuario y simulación apoyan las consultas de
historial y de cuotas. La ausencia de una relación SQL entre
\texttt{usuario\_id} y \texttt{auth\_db.users} es coherente con la separación
de bases, pero permite que un identificador huérfano exista si se eliminara
un usuario por fuera del servicio. Un sistema productivo requeriría una
política explícita para ese caso.

\section{Métodos de amortización}

\subsection{Convenciones}

Sea \(P\) el principal, \(n\) el número de meses, \(r\) la tasa mensual y
\(B_k\) el saldo después de la cuota \(k\). El código interpreta la tasa
recibida como porcentaje anual y calcula
\begin{equation}
r = \frac{\text{tasa anual}}{100\cdot 12}.
\label{eq:tasa}
\end{equation}
Los importes se redondean a dos decimales con redondeo ``away from zero''. La
última cuota ajusta el capital al saldo restante, lo que hace que el saldo
final sea \(0.00\) bajo el algoritmo implementado.

\subsection{Sistema francés}

En el método francés la cuota teórica es constante:
\begin{equation}
C = \frac{P r}{1-(1+r)^{-n}}, \qquad r\ne 0.
\label{eq:frances}
\end{equation}
Para \(r=0\), el código utiliza \(C=P/n\). En cada periodo se calcula
\begin{align}
I_k &= \operatorname{round}(B_{k-1}r),\\
A_k &= \operatorname{round}(C-I_k),\\
B_k &= \operatorname{round}(B_{k-1}-A_k).
\end{align}
La implementación redondea la cuota, el interés, el capital y el saldo en
cada iteración. Por ello, la cuota puede variar en centavos, especialmente
en la última fila, aunque el comportamiento esperado sea aproximadamente
constante.

\subsection{Sistema alemán}

En el método alemán el abono a capital es aproximadamente constante:
\begin{equation}
A = \operatorname{round}\left(\frac{P}{n}\right).
\label{eq:aleman}
\end{equation}
La cuota disminuye porque el interés se calcula sobre un saldo cada vez menor:
\begin{align}
I_k &= \operatorname{round}(B_{k-1}r),\\
C_k &= \operatorname{round}(A+I_k),\\
B_k &= \operatorname{round}(B_{k-1}-A_k).
\end{align}
El último abono también se reemplaza por el saldo restante. Esta estrategia
evita saldos residuales, pero significa que el valor exacto del abono puede
variar en la última cuota por efecto del redondeo.

\subsection{Decisiones discutibles del cálculo}

La tasa anual se convierte a tasa mensual nominal mediante división entre
12; no se documenta una conversión efectiva anual. Esa elección es razonable
para un prototipo, pero debe hacerse explícita al comparar resultados con una
entidad financiera. Además, el proyecto calcula sobre meses uniformes y no
considera fechas de vencimiento, días transcurridos, mora, seguros o
comisiones. Por tanto, una tabla correcta según estas fórmulas no equivale a
una liquidación bancaria.

\section{Infraestructura y configuración}

Docker Compose levanta PostgreSQL 17 Alpine, los tres servicios .NET, el
frontend con Nginx y pgAdmin. PostgreSQL ejecuta los scripts de
\texttt{infrastructure/postgres/init} la primera vez que se crea el volumen.
Los volúmenes \texttt{postgres\_data} y \texttt{pgadmin\_data} conservan
datos y configuración.

\begin{table}[H]
\centering
\caption{Puertos publicados por la composición.}
\label{tab:puertos}
\begin{tabular}{lll}
\hline
\textbf{Componente} & \textbf{Puerto} & \textbf{Uso}\\
\hline
Frontend & 8080 & Aplicación web\\
Auth Service & 5001 & API y Swagger\\
Credit Catalog Service & 5002 & API y Swagger\\
Simulation Service & 5003 & API y Swagger\\
PostgreSQL & 5432 & Desarrollo local\\
pgAdmin & 5050 & Administración de PostgreSQL\\
\hline
\end{tabular}
\end{table}

El arranque documentado es:
\begin{verbatim}
docker compose -f infrastructure/docker-compose.yml up -d --build
\end{verbatim}
Las variables \texttt{VITE\_AUTH\_API\_URL},
\texttt{VITE\_CATALOG\_API\_URL} y
\texttt{VITE\_SIMULATION\_API\_URL} permiten cambiar las URLs del frontend.
El entorno de Compose mostrado contiene credenciales de desarrollo y una
clave JWT explícita. Es una configuración conveniente para una demostración,
pero no debe trasladarse a producción.

\section{Seguridad}

Hay varias medidas positivas: BCrypt para contraseñas, JWT con validación de
firma, emisor, audiencia y expiración, autorización en el controlador de
simulaciones y filtros por usuario en las consultas. También se valida el
método y se acotan monto y plazo en el DTO.

El principal riesgo visible es de configuración. La clave JWT de desarrollo y
las credenciales de PostgreSQL/pgAdmin aparecen en archivos del proyecto o en
la composición. Aunque pueden ser datos deliberadamente didácticos, deben
reemplazarse por secretos gestionados y rotables. También se permite CORS
desde cualquier origen en los servicios. Para producción convendría restringir
orígenes, usar HTTPS, separar secretos por entorno, revisar políticas de
registro de errores, añadir límites de tasa y definir una estrategia para
revocar o renovar tokens.

Guardar el JWT en \texttt{localStorage} simplifica el frontend, pero aumenta
el impacto de una vulnerabilidad XSS. Una alternativa sería una cookie
HttpOnly, Secure y SameSite, acompañada de una defensa CSRF apropiada. Esta es
una mejora arquitectónica, no una funcionalidad actualmente implementada.

\section{Pruebas y validación}

El repositorio no incluye proyectos de pruebas automatizadas ni scripts de
CI visibles. En consecuencia, no se puede afirmar cobertura de código ni
regresión automática. El README sí enumera comprobaciones manuales de
PostgreSQL, catálogo, registro, login, simulación francesa, historial, CSV,
compilación de Vite y levantamiento de Compose. Esas comprobaciones son útiles
como evidencia de desarrollo, pero deberían convertirse en pruebas
reproducibles.

Una estrategia mínima recomendada sería:
\begin{enumerate}[label=\arabic*.]
  \item pruebas unitarias para las fórmulas, tasa cero, redondeo y saldo final;
  \item pruebas de API para credenciales inválidas, métodos inválidos y
  aislamiento entre usuarios;
  \item pruebas de integración con PostgreSQL efímero;
  \item una prueba de frontend para login, ruta privada y exportaciones;
  \item un flujo de CI que ejecute restauración, compilación y pruebas.
\end{enumerate}

\section{Fortalezas, debilidades y recomendaciones}

\begin{table}[H]
\centering
\caption{Evaluación crítica del estado actual.}
\label{tab:evaluacion}
\begin{tabular}{>{\raggedright\arraybackslash}p{0.25\textwidth}
>{\raggedright\arraybackslash}p{0.29\textwidth}
>{\raggedright\arraybackslash}p{0.29\textwidth}}
\hline
\textbf{Área} & \textbf{Fortaleza} & \textbf{Debilidad o mejora}\\
\hline
Arquitectura & Responsabilidades y bases separadas. &
Mayor complejidad operativa y consistencia distribuida.\\
Cálculo & Fórmulas claras y ajuste del saldo final. &
Sin comisiones, fechas reales ni tasa efectiva.\\
Seguridad & BCrypt, JWT y aislamiento por usuario. &
Secretos de desarrollo, CORS abierto y token en almacenamiento local.\\
Persistencia & Snapshot de nombre y tasa histórica. &
Sin integridad referencial entre bases.\\
Frontend & Flujo simple y exportación CSV/PDF. &
Dos formatos mantenidos por capas diferentes.\\
Calidad & Validaciones de entrada y Swagger en desarrollo. &
No hay pruebas automatizadas visibles ni CI comprobable.\\
\hline
\end{tabular}
\end{table}

Las mejoras con mayor retorno serían parametrizar secretos por entorno,
restringir CORS, añadir pruebas de calculadores y API, documentar con precisión
la convención de tasa, y unificar la generación de reportes. Después podrían
incorporarse observabilidad, reintentos y timeouts explícitos para la llamada
al catálogo, versionado de API y migraciones formales en lugar de depender
solo de scripts de inicialización.

\section{Conclusiones}

Credit Simulator cumple una función educativa completa de extremo a extremo:
una interfaz React solicita datos, los microservicios aplican autenticación,
el catálogo entrega una tasa, Simulation Service calcula y guarda una tabla, y
la persona puede consultar o exportar el resultado. La separación de
servicios y bases hace visible una arquitectura escalable, aunque también
introduce decisiones de consistencia y despliegue que no aparecen en una
aplicación monolítica.

La implementación es un punto de partida sólido para una demostración
académica. Sus límites principales son la configuración de desarrollo
expuesta, la política CORS amplia, la falta de pruebas automatizadas visibles
y el alcance financiero simplificado. Estas limitaciones no invalidan el
simulador, pero sí deben comunicarse antes de interpretar sus resultados como
ofertas o liquidaciones reales.

\begin{thebibliography}{9}

\bibitem{aspnet}
Microsoft, ``ASP.NET Core documentation,'' 2024. [En línea].
Disponible: \url{https://learn.microsoft.com/aspnet/core/}

\bibitem{jwt}
M. Jones, J. Bradley y N. Sakimura, ``JSON Web Token (JWT),'' RFC 7519,
Internet Engineering Task Force, 2015. [En línea].
Disponible: \url{https://www.rfc-editor.org/rfc/rfc7519}

\bibitem{postgres}
The PostgreSQL Global Development Group, ``PostgreSQL documentation,'' 2024.
[En línea]. Disponible: \url{https://www.postgresql.org/docs/}

\bibitem{docker}
Docker, ``Docker Compose documentation,'' 2024. [En línea].
Disponible: \url{https://docs.docker.com/compose/}

\bibitem{react}
Meta Open Source, ``React documentation,'' 2024. [En línea].
Disponible: \url{https://react.dev/}

\bibitem{repo}
J. Redroban, ``Credit Simulator,'' repositorio del proyecto, 2026.

\end{thebibliography}

\end{document}
